% Formatage général
	\documentclass[a4paper,11pt]{article}
	\usepackage[T1]{fontenc} % Encodage des caractères spéciaux
	\usepackage{helvet} % Police équivalent Arial
	\usepackage[french]{babel} % Localisation en français
	\frenchbsetup{StandardLists=true} % Pour dire à babel de garder le formatage standard (anglais) des listes à puce
	\usepackage{csquotes} % Pour remplacer automatiquement les guillemets par des guillemets typographique
	\MakeOuterQuote{"} % Le caractère utilisé pour rentrer les guillemets dans le code est le "
	
% Dimensions
	\usepackage[margin=2cm]{geometry} % Marges
	\setlength{\parindent}{0pt} % Largeur de l'alinéa

% Couleurs
	\usepackage{xcolor} % Permet de passer un texte en couleur avec la commande \textcolor{color}{text}
	\definecolor{darkpowderblue}{rgb}{0.0, 0.2, 0.6} % Couleur perso bleue foncée appliquée aux liens URL
	\definecolor{gamboge}{rgb}{0.89, 0.61, 0.06} % Couleur perso orange appliquée aux subsubsections
	\definecolor{ece}{RGB}{0, 122, 123} % Couleur ECE appliquée aux autres titres
	
% Formatage des titres
	\usepackage{titlesec} % Package utilisé pour formater les titres
	\setcounter{secnumdepth}{3} % Numéortation jusqu'au niveau subsubsection
	
	\titleformat{\section} % Formatage des sections
		{\fontsize{18}{18}\bfseries\color{ece}} % Titre en 18 pt, gras, couleur ECE
		{\Roman{section}.} % Numérotation romaine majuscule suivie d'un point
		{1ex} % Un espace entre le numéro et le titre
		{}
	
	\titleformat{\subsection} % Formatage des sous-sections
		{\fontsize{16}{16}\bfseries\color{ece}} % Titre en 16 pt, gras, couleur ECE
		{\Alph{subsection}.} % Numérotation en lettre majuscule suivie d'un point
		{1ex} % Un espace entre le numéro et le titre
		{}
	
	\titleformat{\subsubsection} % Formatage des sous-sous-sections
		{\fontsize{15}{15}\bfseries\color{gamboge}} % Titre en 15 pt, gras, couleur orange
		{(\alph{subsubsection})} % Numérotation en lettre minuscule entourée de parenthèses
		{1ex} % Un espace entre le numéro et le titre
		{}
	
% En-tête et pied de page
	\usepackage{fancyhdr} % Package utilisé pour personnaliser facilement en-tête et pied-de-page
	\renewcommand{\headrulewidth}{0pt} % Épaisseur de la ligne de séparation entre en-tête et corps du rapport
	
	\fancypagestyle{main}{ % En-tête et pied de page pour toutes les autres pages
		\fancyhead[L]{\includegraphics[width=2cm]{images/ece.png}} % En haut à gauche, le logo ECE
		\fancyhead[R]{ING\colorbox{yellow}{[X]} Groupe\colorbox{yellow}{[X]}} % En haut à droite, l'année et le groupe
		\fancyfoot[C]{\thepage} % En bas au centre, le numéro de page
	}
	\pagestyle{main}
		
% Page de titre
	\newcommand{\HRule}[1]{\rule{\linewidth}{#1}} % Ligne de séparation horizontale d'une épaisseur donnée

% Figures
	\usepackage{graphicx} % Permet l'ajout d'images en spécifiant la taille
	\usepackage{subcaption} % Package utilisé pour avoir des sous-figures avec chacune une légende
	\usepackage[labelfont=bf]{caption} % "Figure" et numéro en gras
	
% Liens
	\usepackage[colorlinks=true,linkcolor=black,urlcolor=darkpowderblue]{hyperref}

\begin{document}
	
	\title{
		\begin{figure}[htb]
			\begin{minipage}[t]{.45\textwidth}
				\centering
				\includegraphics[width=6cm]{images/ece.png}
			\end{minipage}
			\hfill
			\begin{minipage}[t]{.45\textwidth}
				\centering
				\raggedleft\vspace{-12mm}\Large{\textbf{ING\colorbox{yellow}{[X]}}\\ Groupe \colorbox{yellow}{[X]}}
			\end{minipage}  
		\end{figure}
		\vspace{2cm}
		\HRule{1.5pt} \\ [0.5cm]
		\LARGE \textbf{\Large{RAPPORT DE TP}\\ [5mm]
		\huge{\textcolor{ece}{TP \colorbox{yellow}{[X]} : \colorbox{yellow}{[Nom du TP]}}}}
		\HRule{1.5pt} \\ [3mm]
		
		\normalsize
		\vspace{1cm}
		\begin{figure}[!h]
			\centering
			\includegraphics[width=\textwidth]{images/elec.png}
		\end{figure}
		\vfill
	}
	
	\author{
		\colorbox{yellow}{André-Marie AMPÈRE} \and \colorbox{yellow}{Alessandro VOLTA}
	}
	
	\date{
		Nous attestons que ce travail est original,\\
		qu’il est le fruit d’un travail commun au binôme et qu’il a été rédigé de manière autonome.\\ [2mm]
		\textbf{\colorbox{yellow}{Ville, le JJ/MM/AAAA}}
	}
	
	\maketitle % Page de titre
	
	\newpage		
	\section{Nom de la partie}
		\subsection{Nom de la section}
	
			\textbf{\textcolor{ece}{T1.}} Réponse à la question T1.
			
			\begin{figure}[!h]
				\centering
				\includegraphics[width=5cm]{images/ece.png}
				\caption{Légende de la Figure}
			\end{figure}
				
			\vspace{3mm}
			\textbf{\textcolor{ece}{E1.}} Réponse à la question E1.
			
			\begin{figure}[!h]
				\centering
				\includegraphics[width=5cm]{images/ece.png}
				\caption{Légende de la Figure}
			\end{figure}
		
	\section{Nom de la partie}
		\subsection{Nom de la section}
	
			\textcolor{gray}{Petits rappels pour une grande réussite :}
			
			\vspace{2mm}
			
			\textcolor{gray}{Pour chaque question, dire :}
			
			\begin{enumerate}
				\item \textcolor{gray}{Formulation du \textbf{problème} : ce que l’on cherche à faire ;}
				\item \textcolor{gray}{Présentation de la \textbf{solution technique} : comment on le fait (bien souvent, algorigramme commenté, ou diagramme de branchement) ;}
				\item \textcolor{gray}{Présentation et description des \textbf{résultats} : Ce que l’on obtient ;}
				\item \textcolor{gray}{\textbf{Validation} des résultats et \textbf{Réponses à la question} (lien entre 1. et 3.) avec un avis critique.}
			\end{enumerate}
			
			\vspace{6mm}
			
			\begin{itemize}
				\item \textcolor{gray}{Toujours mentionner les renvois aux Figures ("Les résultats de la simulation sont présentés sur la Figure 2.1. On y voit…").}
				\item \textcolor{gray}{Pas de code dans les rapports, uniquement des renvois à des annexes si cela est pertinent.}
				\item \textcolor{gray}{Pas de vidéo. Elles ne seront pas regardées et donc pas notées.}
				\item \textcolor{gray}{Lire scrupuleusement l’énoncé. Une mesure à l’oscilloscope n’est pas une capture d’écran du traceur série.}
				\item \textcolor{gray}{Exporter au format .PDF le rapport.}
			\end{itemize}
	
	\newpage
	\section{Annexes}
		Documents volumineux, éventuels codes (\textbf{\textcolor{red}{pas de code dans le rapport}}).

\end{document}